\section{Diagnostic Protocol}
\label{sec:protocol}

Our detection protocol starts from a confound that aggregate AUROC hides. Let an object mention be generated at position $p$, and let $y$ indicate whether CHAIR marks it hallucinated. On LLaVA-1.5-7B / MSCOCO, the hallucination rate rises from $2.2\%$ in the first ten generated tokens to above $60\%$ after position 100. At the same time, many attention statistics drift with decoding position as the model moves from image-conditioned prefix processing toward generated text. Thus position acts as a common cause $p\rightarrow s$ and $p\rightarrow y$ for a detector score $s$: a score can obtain high AUROC by tracking where the object appears in the caption rather than whether the image supports it. The confound is large enough that position alone reaches $0.830$ AUROC.

This confound is not a small calibration nuisance. AUROC is a pairwise ranking metric: it asks whether a random hallucinated mention receives a higher score than a random grounded mention. If hallucinated mentions are late and grounded mentions are early, a position-tracking score will win many cross-position pairs while failing all same-position comparisons. This is why we do not treat overall AUROC as a sufficient hallucination-detection metric. It remains useful for comparability with prior work, but it cannot establish that a detector reads visual evidence.

We therefore report aggregate metrics only as a diagnostic baseline and use four position-controlled metrics as the main test. \textbf{Within-bin AUROC} computes AUROC inside generated-token bins of width 10 and averages valid bins by size. \textbf{Matched-pair AUROC} compares hallucinated mentions only against grounded mentions within $\pm5$ generated tokens. \textbf{Residual AUROC} removes $E[s\mid p]$ before scoring; we use fixed-bin residuals by default and also report linear, cubic, and quantile-bin variants. \textbf{Same-object matching} further restricts matched pairs to the same generated object word, addressing category-frequency confounds. A genuinely grounded detector should retain discrimination under these controls; a positional proxy should collapse toward chance.

Our mitigation protocol mirrors the same logic. Standard POPE accuracy or F1 can improve when a method shifts the yes/no answer prior, so we report yes rate, TPR, FPR, MCC, and $\Delta\mathrm{TPR}-\Delta\mathrm{FPR}$ rather than treating recall gains alone as mitigation. CHAIR captioning is evaluated with CHAIR$_i$, CHAIR$_s$, mean caption length, object mentions, and hallucinated-object mentions, because a method can lower hallucinations by producing shorter or less object-rich captions. Finally, we add an attention-shift audit on adversarial POPE: for each intervention we record pre/post visual, prefix, and sink attention mass in active layers and compare true positives against false positives. This audit tests whether the intended routing change is selective for verified object presence.

The associated-evidence audit adds a harder verification slice. Instead of asking only whether the target object is absent, it asks whether a method remains reliable when a related object or scene context is present. These semantic-neighbor negatives are the cases in which attention can look most persuasive while still failing the target-object claim. A detector should not mark them as grounded simply because the attention map is coherent; a mitigation method should not increase their yes rate simply because the image contains plausible associated evidence.

The protocol is intentionally diagnostic rather than leaderboard-oriented. It follows the spirit of the contrastive-decoding critique in \citet{yin2025mirage}: when a benchmark reward can be obtained by a simpler spurious mechanism, the evaluation must separate intended causal effect from shortcut behavior. For detection, the shortcut is generation position. For mitigation, the shortcut is answer-prior or caption-style movement. In both cases, the question is the same: does the method change object-level evidence verification, or merely a correlated surface statistic?

\paragraph{Setup.} All main experiments use LLaVA-1.5-7B \cite{liu2023llava}. Detection uses MSCOCO val2014 \cite{lin2014coco} captions generated with the standard LLaVA prompt, then labels $16{,}426$ object mentions ($4{,}009$ hallucinated) with CHAIR \cite{rohrbach2018chair}. Detection baselines include NLL, entropy, IC \cite{jiang2024interpreting}, GLSim \cite{park2025glsim}, SVAR, PAS, Beyond-ADS/CGC, SinkDetect, and OWLv2-TDEV region evidence. Mitigation uses COCO POPE random/popular/adversarial splits and CHAIR caption generation over the same $5{,}000$ COCO images. We evaluate PAI-attention-only, ClearSight VAF, Visual Attention Sink, VCD-greedy, NoLan-compatible, and TDEV direct, gate, and hybrid gate-plus-rescue rules under one HuggingFace LLaVA stack; PAI is explicitly the attention component, excluding its contrastive/CFG branch, VCD-greedy is a deterministic port of the contrastive logit rule, and NoLan-compatible is a deterministic local port of dynamic language-prior suppression rather than an official reproduction.
